\section{Phân tích lỗi và thảo luận}
\label{sec:discussion}

\subsection{Tổng quan các nguồn lỗi}

Kết quả thực nghiệm cho thấy lỗi của pipeline không tập trung tại một thành
phần duy nhất mà có thể lan truyền qua nhiều bước. Một lỗi retrieval có thể
làm mất gold event hoặc gold rule; từ đó gold path không xuất hiện trong
candidate set, verifier xác minh một path khác và answer generator tạo câu
trả lời nhất quán với path sai đó.

Ngược lại, ngay cả khi gold path đã xuất hiện trong candidate set, hệ thống
vẫn có thể chọn sai primary path. Trong trường hợp này, retrieval coverage
có thể cao nhưng Top-1 Exact Path, answer và citation vẫn thấp. Các nhóm lỗi
chính được tóm tắt trong Bảng~\ref{tab:error-summary}.

\begin{table}[t]
    \centering
    \caption{Các nguồn lỗi chính được quan sát trong thực nghiệm.}
    \label{tab:error-summary}
    \begin{tabular}{p{0.28\linewidth}p{0.22\linewidth}p{0.38\linewidth}}
        \hline
        \textbf{Nguồn lỗi} & \textbf{Tín hiệu định lượng} &
        \textbf{Ảnh hưởng} \\
        \hline
        Path retrieval &
        Error rate \(28{,}80\%\) &
        Gold path không xuất hiện trong candidate set. \\

        Path ranking &
        Error rate \(36{,}00\%\) &
        Gold path đã được truy hồi nhưng không được chọn làm primary path. \\

        Reverse reasoning &
        Answer F1 \(32{,}67\%\) &
        Sai seed hoặc sai predecessor dẫn đến sai path, decision và answer. \\

        Verification abstention &
        \(20\) mẫu supported bị dự đoán uncertain &
        Giảm accuracy nhưng tránh tạo quyết định khẳng định sai. \\

        Citation selection &
        Citation F1 \(62{,}86\%\) &
        Thiếu gold article hoặc bổ sung article từ path cạnh tranh. \\

        Metric cho cấu trúc phi tuyến &
        Exact path bằng \(0\%\) ở branch/convergence &
        Metric tuyến tính không phản ánh đầy đủ chất lượng câu trả lời. \\
        \hline
    \end{tabular}
\end{table}

\subsection{Khoảng cách giữa path retrieval và path ranking}

Một trong những kết quả đáng chú ý nhất là khoảng cách giữa Top-1 Exact Path
và Oracle Exact Path. Top-1 chỉ đạt \(35{,}20\%\), trong khi Oracle đạt
\(71{,}20\%\). Như vậy, trong một tỷ lệ lớn câu hỏi, retriever đã tạo được
một candidate path phù hợp với gold nhưng scoring function chưa xếp path đó
ở vị trí đầu tiên.

Có ba nguyên nhân chính có thể giải thích hiện tượng này.

\paragraph{Nhiều path có chung prefix hoặc mediator.}

Một event trung gian có thể dẫn đến nhiều outcome khác nhau thông qua các
rule thuộc nhiều điều luật. Ví dụ, các event như ``tái phạm nguy hiểm'' hoặc
``phạm tội có tổ chức'' có thể liên kết với nhiều khung hình phạt. Các path
này có chung seed hoặc mediator nên đạt semantic và graph score gần nhau.
Nếu truy vấn chứa thông tin phân biệt outcome nhưng thông tin đó chưa được
đưa đủ mạnh vào scoring function, hệ thống có thể chọn một path hợp lệ về
cấu trúc nhưng sai về hệ quả cụ thể.

\paragraph{Outcome có nội dung ngữ nghĩa gần nhau.}

Nhiều effect event cùng mô tả hình phạt tù, phạt tiền hoặc trách nhiệm hình
sự nhưng khác nhau về mức hình phạt và điều luật. Embedding có thể đánh giá
các event này là gần nhau trong không gian vector. Một khác biệt nhỏ về con
số hoặc phạm vi hình phạt có ý nghĩa pháp lý lớn nhưng không nhất thiết tạo
ra khoảng cách embedding đủ lớn.

\paragraph{Path score chưa sử dụng đầy đủ answer intent.}

Retriever chạy trước khi Step 5 xác định answer intent. Mặc dù truy vấn đã
được sử dụng trong semantic retrieval, scoring function chưa áp đặt ràng
buộc rõ ràng rằng reverse question phải ưu tiên path có outcome khớp truy
vấn và trả về predecessor, hoặc branch question phải giữ nhiều path có chung
prefix. Vì vậy, path ranking chủ yếu dựa trên relevance và graph score thay
vì vai trò mà từng event cần đóng trong câu trả lời.

Kết quả Oracle cho thấy chỉ cải thiện candidate generation là chưa đủ. Một
hướng ưu tiên cao hơn là xây dựng query-aware path reranker, trong đó query
intent được xác định trước hoặc đồng thời với retrieval. Reranker có thể sử
dụng các đặc trưng:

\begin{itemize}
    \item mức khớp giữa seed event và phần nguyên nhân của truy vấn;
    \item mức khớp giữa outcome event và phần kết quả của truy vấn;
    \item vai trò kỳ vọng của event được hỏi;
    \item mức khớp với rule condition và rule effect đầy đủ;
    \item article hoặc mức hình phạt xuất hiện trong câu hỏi;
    \item tính liên tục của toàn bộ path;
    \item độ đặc hiệu của event thay vì chỉ semantic similarity tổng quát.
\end{itemize}

\subsection{Lỗi path retrieval}

Path Retrieval Error Rate đạt \(28{,}80\%\), nghĩa là gần ba phần mười số mẫu
không có gold path hoàn chỉnh trong candidate set. Lỗi này không thể được
khắc phục chỉ bằng cách thay đổi Step 4 hoặc Step 5.

Một số nguyên nhân gồm:

\begin{itemize}
    \item seed event đúng không nằm trong top-\(k\);
    \item semantic score thấp do cách diễn đạt trong câu hỏi khác event name;
    \item candidate pool bị chiếm bởi nhiều event có nội dung hình phạt gần
    nhau;
    \item graph traversal bắt đầu từ một event đúng về chủ đề nhưng sai vai
    trò;
    \item giới hạn số path trên mỗi seed loại bỏ gold path trước bước
    reranking;
    \item gold rule không xuất hiện đủ trong candidate rule pool.
\end{itemize}

Rule Recall@5 và Event Recall@5 lần lượt đạt \(69{,}93\%\) và \(74{,}66\%\),
cho thấy retriever thường lấy được một phần path nhưng chưa luôn lấy đủ mọi
hop. Điều này phù hợp với Top-1 Edge F1 \(54{,}13\%\), cao hơn Exact Path
nhưng vẫn thấp hơn Oracle Event F1.

Các cải tiến retrieval có thể bao gồm:

\begin{enumerate}
    \item tách query thành biểu diễn nguyên nhân, mediator và outcome trước
    khi tìm seed;

    \item thực hiện retrieval riêng cho condition text và effect text;

    \item kết hợp dense score với sparse lexical matching cho số điều,
    khoảng tiền và mức hình phạt;

    \item tăng tính đa dạng của seed event thay vì lấy nhiều event gần nhau
    về ngữ nghĩa;

    \item thực hiện constrained backward search cho reverse reasoning;

    \item dùng rule-aware reranker trên toàn bộ path thay vì xếp từng event
    độc lập.
\end{enumerate}

\subsection{Reverse reasoning là điểm yếu chính}

Reverse reasoning là nhóm có hiệu năng thấp nhất ở hầu hết metric:

\begin{itemize}
    \item Rule Recall@5: \(46{,}67\%\);
    \item Event Recall@5: \(42{,}22\%\);
    \item Top-1 Exact Path: \(8{,}89\%\);
    \item Oracle Exact Path: \(26{,}67\%\);
    \item Verification Accuracy: \(60{,}00\%\);
    \item Answer Token F1: \(32{,}67\%\);
    \item Citation F1: \(36{,}32\%\).
\end{itemize}

Trong forward reasoning, truy vấn thường nêu rõ condition event và yêu cầu
tìm outcome. Condition event có thể được dùng trực tiếp làm seed, sau đó
graph được mở rộng theo hướng \texttt{CAUSES}. Trong reverse reasoning, truy
vấn thường bắt đầu từ outcome và yêu cầu tìm nguyên nhân. Một outcome như
``chịu trách nhiệm hình sự'', ``được giảm nhẹ'' hoặc một mức hình phạt có thể
có nhiều predecessor thuộc các điều luật khác nhau. Do đó, bài toán ngược có
mức nhập nhằng cao hơn.

Mặc dù retriever được cấu hình \texttt{direction=both}, việc cho phép duyệt
cạnh ngược chưa đủ để giải quyết vấn đề. Hệ thống còn phải xác định rằng
outcome được nhắc đến trong truy vấn nên được cố định làm điểm cuối, sau đó
xếp hạng các predecessor dựa trên toàn bộ nội dung câu hỏi.

Trong \(45\) reverse question, Verification Accuracy \(60\%\) tương ứng với
\(27\) quyết định đúng và \(18\) mẫu bị chuyển thành
\texttt{UNCERTAIN}. Hai mẫu \texttt{SUPPORTED} còn lại bị chuyển thành
\texttt{UNCERTAIN} thuộc nhóm yes/no positive, tạo thành tổng \(20\) lỗi xác
minh trong confusion matrix.

Answer generator đã hỗ trợ intent \texttt{INITIAL\_CAUSE}, nhưng nó chỉ có
thể trả nguyên nhân đầu tiên của primary path được chọn. Nếu upstream
retriever hoặc verifier chọn sai path, việc định dạng đúng answer intent
không thể khôi phục gold cause. Kết quả này cho thấy lỗi reverse reasoning
chủ yếu nằm trước Step 5.

Một cải tiến trực tiếp là thực hiện backward retrieval có ràng buộc:

\begin{equation}
\widehat{p}
=
\arg\max_{p}
S(p \mid q)
\quad
\text{sao cho}
\quad
\mathrm{outcome}(p)
=
v_q^{\mathrm{outcome}},
\end{equation}

trong đó \(v_q^{\mathrm{outcome}}\) là outcome event được ánh xạ từ truy vấn.
Sau khi outcome được cố định, hệ thống mới xếp hạng các path đi ngược về seed
event.

\subsection{Hành vi abstention của verifier}

Confusion matrix cho thấy verifier không nhầm lẫn trực tiếp giữa hai lớp có
gold support:

\begin{itemize}
    \item không có gold \texttt{SUPPORTED} nào bị dự đoán thành
    \texttt{REJECT\_DIRECT\_CLAIM};

    \item không có gold \texttt{REJECT\_DIRECT\_CLAIM} nào bị dự đoán thành
    \texttt{SUPPORTED};

    \item toàn bộ \(20\) lỗi là gold \texttt{SUPPORTED} bị chuyển thành
    \texttt{UNCERTAIN}.
\end{itemize}

Hành vi này tạo ra precision \(100\%\) cho cả lớp \texttt{SUPPORTED} và
\texttt{REJECT\_DIRECT\_CLAIM}, nhưng làm giảm recall của lớp
\texttt{SUPPORTED} xuống \(91{,}30\%\).

Trong một hệ thống pháp luật, việc chuyển sang \texttt{UNCERTAIN} khi thiếu
evidence có thể an toàn hơn việc đưa ra một kết luận khẳng định sai. Tuy
nhiên, benchmark hiện không chứa gold \texttt{UNCERTAIN}, nên chưa thể đánh
giá liệu threshold abstention có được hiệu chỉnh phù hợp hay không.

Một benchmark tương lai nên chứa ít nhất ba loại mẫu:

\begin{enumerate}
    \item claim được hỗ trợ;
    \item claim bị bác bỏ;
    \item claim thực sự không đủ căn cứ.
\end{enumerate}

Khi đó, có thể đánh giá coverage--risk trade-off thay vì chỉ báo cáo
accuracy và Macro-F1.

\subsection{Hiệu quả của direct-claim detection}

Tất cả \(20\) yes/no counterexample đều được phân loại đúng. Đây là kết quả
tích cực của query-aware claim detection: hệ thống nhận ra rằng câu hỏi yêu
cầu kiểm tra quan hệ trực tiếp và không xem một đường hai bước là bằng chứng
cho cạnh trực tiếp.

Tuy nhiên, kết quả này chủ yếu đến từ việc kết hợp:

\begin{itemize}
    \item direct indicator trong query;
    \item cấu trúc cạnh của graph;
    \item sự phân biệt giữa direct edge và indirect path.
\end{itemize}

Nó không nhất thiết yêu cầu hard intervention. Vì vậy, không nên sử dụng độ
chính xác \(100\%\) trên subset này để kết luận LegalSCM có ưu thế về
counterfactual reasoning.

\subsection{Chất lượng câu trả lời và độ dài đầu ra}

Answer Token Recall đạt \(75{,}40\%\), cao hơn Token Precision
\(65{,}34\%\). Câu trả lời thường bao phủ nội dung gold nhưng bổ sung:

\begin{itemize}
    \item câu dẫn như ``Hệ quả cuối cùng là'';
    \item toàn bộ causal chain;
    \item giải thích quyết định;
    \item citation marker \([E_i]\).
\end{itemize}

Các thành phần bổ sung tăng khả năng giải thích nhưng làm giảm token
precision. Prediction có trung bình \(48{,}18\) token, trong khi gold answer
có trung bình \(45{,}07\) token.

Normalized Exact Match bằng \(0\%\) cho thấy exact string matching không phù
hợp với dạng câu trả lời có giải thích. Token F1 và ROUGE-L F1 phản ánh chất
lượng tốt hơn, nhưng vẫn chưa đánh giá đầy đủ tính đúng đắn pháp lý.

Một hướng cải tiến là tách đầu ra thành hai trường:

\begin{itemize}
    \item \texttt{short\_answer}: câu trả lời trực tiếp, ngắn gọn;
    \item \texttt{explanation}: causal chain và căn cứ.
\end{itemize}

Khi đó, short answer có thể được đánh giá bằng exact match hoặc token F1,
còn explanation được đánh giá riêng bằng path grounding và citation.

\subsection{Ảnh hưởng của event canonicalization}

Một số event name trong graph chứa nhiều cách diễn đạt được nối bằng ký hiệu
\texttt{||}. Ví dụ, một node có thể có tên:

\begin{quote}
``Có thể được miễn phần hình phạt còn lại
\texttt{||}
Được miễn phần hình phạt còn lại''.
\end{quote}

Việc đưa toàn bộ chuỗi canonicalized name vào câu trả lời làm tăng số token
không xuất hiện trong gold answer và giảm tính tự nhiên. Nó cũng có thể làm
semantic retrieval ưu tiên các event có nhiều biến thể ngôn ngữ hơn.

Cần tách rõ:

\begin{itemize}
    \item canonical event ID;
    \item danh sách alias;
    \item preferred display name.
\end{itemize}

Retriever có thể sử dụng toàn bộ alias, trong khi answer generator chỉ sử
dụng preferred display name.

\subsection{Phân tích citation}

Citation F1 đạt \(62{,}86\%\), thấp hơn Answer Token F1. Hai nguồn lỗi chính
là citation miss và extra citation.

\paragraph{Citation miss.}

Nếu retriever chọn đúng event nhưng thiếu một rule trên gold path, câu trả
lời có thể vẫn chứa phần lớn nội dung gold nhưng thiếu article tương ứng.
Điều này làm answer F1 cao hơn citation recall.

\paragraph{Extra citation.}

Khi nhiều rule cùng hỗ trợ một causal edge hoặc nhiều path chia sẻ outcome,
Step 5 có thể giữ thêm evidence đã được xác minh nhưng không thuộc gold
citation set. Nếu citation từ evidence bổ sung xuất hiện trong câu trả lời,
citation precision giảm.

Việc khử trùng theo causal edge đã hạn chế citation lặp lại, nhưng chưa giải
quyết hoàn toàn trường hợp nhiều điều luật cùng dẫn đến event tương tự. Một
cải tiến là chỉ gắn citation theo từng assertion cụ thể trong answer thay vì
thêm tất cả citation ở cuối câu.

Ví dụ, với chain explanation:

\[
A \xrightarrow{r_1} B \xrightarrow{r_2} C,
\]

câu trả lời nên gắn:

\begin{quote}
Bước 1: \(A\) dẫn đến \(B\) [Điều \(a_1\)].
Bước 2: \(B\) dẫn đến \(C\) [Điều \(a_2\)].
\end{quote}

Cách này làm rõ article nào hỗ trợ hop nào và tạo điều kiện đánh giá
citation grounding ở mức assertion.

\subsection{Branch và convergence}

Branch reasoning và convergence reasoning đặt ra vấn đề cho metric exact
path. Một branch có thể được biểu diễn bởi tập path:

\[
\mathcal{P}_{\mathrm{branch}}
=
\{
A \rightarrow B \rightarrow C_1,\,
A \rightarrow B \rightarrow C_2,\,
\ldots
\},
\]

trong khi convergence có dạng:

\[
\mathcal{P}_{\mathrm{convergence}}
=
\{
A_1 \rightarrow B \rightarrow C,\,
A_2 \rightarrow B \rightarrow C,\,
\ldots
\}.
\]

Không có một chuỗi tuyến tính duy nhất đại diện đầy đủ cho các cấu trúc này.
Do đó, Top-1 và Oracle Exact Path có thể bằng \(0\%\) ngay cả khi answer
generator liệt kê đúng nhiều nhánh hoặc nhiều điều kiện hội tụ.

Kết quả thực nghiệm minh họa rõ vấn đề:

\begin{itemize}
    \item branch reasoning có Top-1 Exact Path \(0\%\) nhưng Answer Token F1
    \(69{,}33\%\);

    \item convergence reasoning có Top-1 Exact Path \(0\%\) nhưng Answer
    Token F1 \(84{,}32\%\).
\end{itemize}

Một evaluator phù hợp hơn nên sử dụng:

\begin{itemize}
    \item Path-set Precision, Recall và F1;
    \item branch outcome coverage;
    \item convergence input coverage;
    \item subgraph edge F1;
    \item graph edit distance;
    \item assertion-level citation accuracy.
\end{itemize}

Vì chưa có các metric này, kết quả path của branch và convergence chỉ được
xem là diagnostic.

\subsection{Ảnh hưởng của thành phần dữ liệu lên difficulty}

Nhóm hard đạt Answer Token F1 \(81{,}25\%\), cao hơn nhóm medium
\(63{,}90\%\). Verification Accuracy của nhóm hard cũng đạt \(100\%\), trong
khi nhóm medium đạt \(88{,}95\%\).

Kết quả này không phù hợp với cách diễn giải thông thường rằng mẫu hard phải
khó hơn mẫu medium. Nguyên nhân là difficulty label chịu ảnh hưởng của
question type. Nhóm hard chứa nhiều chain, counterexample và cấu trúc có gold
answer rõ ràng theo template. Nhóm medium lại chứa nhiều reverse question,
là nhóm mà retriever hiện xử lý kém.

Do đó, difficulty hiện tại chưa phải một thang đo đã được hiệu chỉnh độc lập.
Nghiên cứu không sử dụng kết quả này để kết luận hệ thống hoạt động tốt hơn
trên câu hỏi khó. Trong tương lai, difficulty nên được xác định từ các yếu
tố có thể đo được, chẳng hạn:

\begin{itemize}
    \item số candidate path cạnh tranh;
    \item bậc vào và bậc ra của mediator;
    \item độ tương đồng giữa các outcome;
    \item số rule hỗ trợ;
    \item số hop;
    \item mức lexical overlap giữa query và gold event;
    \item mức đồng thuận của chuyên gia.
\end{itemize}

\subsection{Giải thích null quality ablation}

Kết quả paired cho thấy structural SCM và path ablation đạt quality metric
giống hệt nhau. Có một số nguyên nhân hợp lý cho kết quả này.

\paragraph{Retrieval không thay đổi.}

Hai engine được áp dụng sau Step 3. Vì vậy, rule, event và candidate path ban
đầu giống nhau.

\paragraph{Phần lớn factual path đơn giản.}

Corpus hiện chủ yếu chứa rule có một điều kiện dương và một hệ quả dương.
Trên các path đơn giản này, fixed-point inference và reachability thường dẫn
đến cùng kết luận supported.

\paragraph{Counterfactual subset không nhạy với intervention semantics.}

Hai mươi counterfactual sample chủ yếu kiểm tra cạnh trực tiếp. Cả structural
SCM và path ablation có thể bác bỏ claim dựa trên cấu trúc
\(A \rightarrow B \rightarrow C\) mà không cần tạo ra khác biệt về
counterfactual outcome.

\paragraph{Thiếu cơ chế thay thế phức tạp.}

Benchmark chưa bao phủ đầy đủ các tình huống trong đó xóa node và hard
intervention tạo kết quả khác nhau, chẳng hạn:

\begin{itemize}
    \item mediator có nhiều cơ chế đầu vào;
    \item intervention làm kích hoạt một rule thay thế;
    \item outcome có nhiều cơ chế cấu trúc cạnh tranh;
    \item rule có exception hoặc negation;
    \item nhiều mediator bị can thiệp đồng thời;
    \item xung đột làm outcome trở thành \(\mathrm{UNKNOWN}\).
\end{itemize}

\paragraph{Step 5 nhận cùng quyết định và primary path.}

Khi hai verifier trả cùng final decision, primary path và evidence được giữ,
answer generator tạo cùng answer và citation. Vì vậy, answer metric không
thể phân biệt hai engine.

Kết quả này không chứng minh hai phương pháp tương đương về mặt lý thuyết.
Hard intervention thay thế mechanism của biến, trong khi node deletion chỉ
thay đổi topology và reachability của graph
\cite{pearl2009causality}. Null result chỉ cho thấy benchmark hiện tại không
làm lộ ra khác biệt đó.

\subsection{Đánh đổi giữa semantics và runtime}

Structural SCM cần trung bình \(3{,}890\) giây trên mỗi mẫu, trong khi path
ablation cần \(1{,}205\) giây. Structural SCM chậm hơn khoảng \(3{,}23\) lần.

Chi phí bổ sung đến từ:

\begin{itemize}
    \item khởi tạo hoặc truy cập LegalSCM;
    \item xây dựng factual context;
    \item fixed-point inference;
    \item kiểm tra rule coverage;
    \item thực hiện intervention cho từng mediator;
    \item suy luận lại counterfactual world;
    \item tạo trace chi tiết;
    \item tính node-deletion baseline diagnostic.
\end{itemize}

Đổi lại, structural SCM lưu được thông tin mà path ablation không cung cấp:

\begin{itemize}
    \item factual state của mediator và outcome;
    \item intervention assignment;
    \item counterfactual outcome;
    \item outcome có thay đổi hay không;
    \item rule bị vô hiệu hóa;
    \item rule mới được kích hoạt;
    \item conflicting event;
    \item số vòng lặp fixed point.
\end{itemize}

Các trường này có thể hữu ích cho auditability và phân tích lỗi. Tuy nhiên,
nghiên cứu hiện chưa có expert evaluation để định lượng giá trị của trace.
Vì vậy, không nên chuyển lợi ích về khả năng diễn giải thành một tuyên bố về
độ chính xác.

\subsection{Yêu cầu đối với benchmark phản thực tế tương lai}

Để đánh giá đúng LegalSCM, benchmark tương lai cần gán ground truth không chỉ
cho final decision mà còn cho world state và intervention trace. Một
counterfactual sample nên chứa:

\begin{itemize}
    \item factual context;
    \item mediator được can thiệp;
    \item factual mediator state;
    \item factual outcome;
    \item intervention assignment;
    \item counterfactual outcome;
    \item outcome changed flag;
    \item các mechanism bị vô hiệu hóa;
    \item các mechanism thay thế;
    \item giải thích được chuyên gia xác nhận.
\end{itemize}

Benchmark cũng cần chứa các cặp trường hợp mà node deletion và hard
intervention cho kết quả khác nhau. Chỉ khi đó paired ablation mới có khả
năng kiểm tra giá trị gia tăng của structural semantics.

\subsection{Hàm ý đối với triển khai thực tế}

Kết quả hiện tại cho thấy pipeline có thể tạo câu trả lời grounded trên toàn
bộ benchmark và phân loại đúng các direct-edge negative claim. Tuy nhiên,
Top-1 Exact Path \(35{,}20\%\), Citation F1 \(62{,}86\%\) và hiệu năng thấp
trên reverse reasoning chưa đủ để hỗ trợ triển khai như một hệ thống tư vấn
pháp lý tự động.

Trong một ứng dụng thực tế, output nên được trình bày như một gợi ý có căn
cứ và cần cho phép người dùng:

\begin{itemize}
    \item xem rule và article nguồn;
    \item kiểm tra từng hop của causal path;
    \item xem các supporting path khác;
    \item nhận biết khi hệ thống trả \texttt{UNCERTAIN};
    \item phân biệt factual inference với counterfactual intervention;
    \item biết phiên bản corpus và thời điểm cập nhật.
\end{itemize}

Hệ thống không nên ẩn uncertainty hoặc chỉ hiển thị final answer mà không
cho phép kiểm tra evidence.

\subsection{Ưu tiên cải tiến}

Từ kết quả thực nghiệm, các hướng cải tiến được ưu tiên theo thứ tự sau:

\begin{enumerate}
    \item \textbf{Query-aware path reranking.}
    Đưa answer intent và event role vào scoring trước khi chọn primary path.

    \item \textbf{Backward retrieval cho reverse reasoning.}
    Cố định outcome được nhắc trong truy vấn và tìm predecessor có điều kiện.

    \item \textbf{Hybrid lexical--dense retrieval.}
    Bổ sung sparse matching cho số điều, mức tiền và khung hình phạt.

    \item \textbf{Preferred event display name.}
    Tách canonical ID, alias và tên hiển thị để loại chuỗi
    \texttt{||} khỏi câu trả lời.

    \item \textbf{Assertion-level citation.}
    Gắn từng article với hop hoặc assertion cụ thể.

    \item \textbf{Subgraph evaluation.}
    Bổ sung metric path-set cho branch và convergence.

    \item \textbf{Intervention-sensitive benchmark.}
    Xây dựng gold factual/counterfactual world có thẩm định chuyên gia.

    \item \textbf{Mở rộng LegalSCM.}
    Hỗ trợ nhiều điều kiện, negation, exception, conflict và rule priority.
\end{enumerate}

Các cải tiến này nên được triển khai trong một iteration mới và đánh giá bằng
một paired experiment mới. Chúng không được trộn với kết quả v4.1/v5.1 đã
đóng băng trong nghiên cứu hiện tại.
